% Options for packages loaded elsewhere
\PassOptionsToPackage{unicode}{hyperref}
\PassOptionsToPackage{hyphens}{url}
%
\documentclass[UTF8]{ctexart}
\usepackage{amsmath,amssymb}
\usepackage{lmodern}
\usepackage{iftex}
\ifPDFTeX
  \usepackage[T1]{fontenc}
  \usepackage[utf8]{inputenc}
  \usepackage{textcomp} % provide euro and other symbols
\else % if luatex or xetex
  \usepackage{unicode-math}
  \defaultfontfeatures{Scale=MatchLowercase}
  \defaultfontfeatures[\rmfamily]{Ligatures=TeX,Scale=1}
  \usepackage{fontspec}
  \usepackage{xeCJK}
  \setmainfont{Times New Roman}
  \setmonofont{DejaVu Sans Mono}
  \setCJKmainfont{SimSun}
\fi
% Use upquote if available, for straight quotes in verbatim environments
\IfFileExists{upquote.sty}{\usepackage{upquote}}{}
\IfFileExists{microtype.sty}{% use microtype if available
  \usepackage[]{microtype}
  \UseMicrotypeSet[protrusion]{basicmath} % disable protrusion for tt fonts
}{}
\makeatletter
\@ifundefined{KOMAClassName}{% if non-KOMA class
  \IfFileExists{parskip.sty}{%
    \usepackage{parskip}
  }{% else
    \setlength{\parindent}{0pt}
    \setlength{\parskip}{6pt plus 2pt minus 1pt}}
}{% if KOMA class
  \KOMAoptions{parskip=half}}
\makeatother
\usepackage{xcolor}
\usepackage[top=2.5cm,bottom=2.5cm,left=3cm,right=3cm]{geometry}
\IfFileExists{xurl.sty}{\usepackage{xurl}}{} % add URL line breaks if available
\IfFileExists{bookmark.sty}{\usepackage{bookmark}}{\usepackage{hyperref}}
\hypersetup{
  hidelinks,
  pdfcreator={LaTeX via pandoc}}
\urlstyle{same} % disable monospaced font for URLs
\usepackage{color}
\usepackage{fancyvrb}
\newcommand{\VerbBar}{|}
\newcommand{\VERB}{\Verb[commandchars=\\\{\}]}
\DefineVerbatimEnvironment{Highlighting}{Verbatim}{commandchars=\\\{\}}
% Add ',fontsize=\small' for more characters per line
\newenvironment{Shaded}{}{}
\newcommand{\AlertTok}[1]{\textcolor[rgb]{1.00,0.00,0.00}{\textbf{#1}}}
\newcommand{\AnnotationTok}[1]{\textcolor[rgb]{0.38,0.63,0.69}{\textbf{\textit{#1}}}}
\newcommand{\AttributeTok}[1]{\textcolor[rgb]{0.49,0.56,0.16}{#1}}
\newcommand{\BaseNTok}[1]{\textcolor[rgb]{0.25,0.63,0.44}{#1}}
\newcommand{\BuiltInTok}[1]{#1}
\newcommand{\CharTok}[1]{\textcolor[rgb]{0.25,0.44,0.63}{#1}}
\newcommand{\CommentTok}[1]{\textcolor[rgb]{0.38,0.63,0.69}{\textit{#1}}}
\newcommand{\CommentVarTok}[1]{\textcolor[rgb]{0.38,0.63,0.69}{\textbf{\textit{#1}}}}
\newcommand{\ConstantTok}[1]{\textcolor[rgb]{0.53,0.00,0.00}{#1}}
\newcommand{\ControlFlowTok}[1]{\textcolor[rgb]{0.00,0.44,0.13}{\textbf{#1}}}
\newcommand{\DataTypeTok}[1]{\textcolor[rgb]{0.56,0.13,0.00}{#1}}
\newcommand{\DecValTok}[1]{\textcolor[rgb]{0.25,0.63,0.44}{#1}}
\newcommand{\DocumentationTok}[1]{\textcolor[rgb]{0.73,0.13,0.13}{\textit{#1}}}
\newcommand{\ErrorTok}[1]{\textcolor[rgb]{1.00,0.00,0.00}{\textbf{#1}}}
\newcommand{\ExtensionTok}[1]{#1}
\newcommand{\FloatTok}[1]{\textcolor[rgb]{0.25,0.63,0.44}{#1}}
\newcommand{\FunctionTok}[1]{\textcolor[rgb]{0.02,0.16,0.49}{#1}}
\newcommand{\ImportTok}[1]{#1}
\newcommand{\InformationTok}[1]{\textcolor[rgb]{0.38,0.63,0.69}{\textbf{\textit{#1}}}}
\newcommand{\KeywordTok}[1]{\textcolor[rgb]{0.00,0.44,0.13}{\textbf{#1}}}
\newcommand{\NormalTok}[1]{#1}
\newcommand{\OperatorTok}[1]{\textcolor[rgb]{0.40,0.40,0.40}{#1}}
\newcommand{\OtherTok}[1]{\textcolor[rgb]{0.00,0.44,0.13}{#1}}
\newcommand{\PreprocessorTok}[1]{\textcolor[rgb]{0.74,0.48,0.00}{#1}}
\newcommand{\RegionMarkerTok}[1]{#1}
\newcommand{\SpecialCharTok}[1]{\textcolor[rgb]{0.25,0.44,0.63}{#1}}
\newcommand{\SpecialStringTok}[1]{\textcolor[rgb]{0.73,0.40,0.53}{#1}}
\newcommand{\StringTok}[1]{\textcolor[rgb]{0.25,0.44,0.63}{#1}}
\newcommand{\VariableTok}[1]{\textcolor[rgb]{0.10,0.09,0.49}{#1}}
\newcommand{\VerbatimStringTok}[1]{\textcolor[rgb]{0.25,0.44,0.63}{#1}}
\newcommand{\WarningTok}[1]{\textcolor[rgb]{0.38,0.63,0.69}{\textbf{\textit{#1}}}}
\setlength{\emergencystretch}{3em} % prevent overfull lines
\providecommand{\tightlist}{%
  \setlength{\itemsep}{0pt}\setlength{\parskip}{0pt}}
\setcounter{secnumdepth}{0}
\ifLuaTeX
  \usepackage{selnolig}  % disable illegal ligatures
\fi

\author{}
\date{}



\begin{document}

% {
% \setcounter{tocdepth}{3}
% \tableofcontents
% }
\hypertarget{saads-db-ux6a21ux5757ux8be6ux7ec6ux8bbeux8ba1ux65b9ux6848}{%
\section{\texorpdfstring{SAADS \texttt{db/}
模块详细设计方案}{SAADS db/ 模块详细设计方案}}\label{saads-db-ux6a21ux5757ux8be6ux7ec6ux8bbeux8ba1ux65b9ux6848}}

\hypertarget{ux8bbeux8ba1ux76eeux6807}{%
\subsection{1. 设计目标}\label{ux8bbeux8ba1ux76eeux6807}}

本方案面向 \textbf{WP1-1 情报采集智能体} 的 Python
实现，目标不是简单``连库 + 执行 SQL''，而是把 PostgreSQL 中已经定义好的
\texttt{wp11} schema，稳定映射为一套可维护、可测试、可被多个 Agent
复用的数据库访问层。

本设计直接对齐当前数据库对象：

\begin{itemize}
\tightlist
\item
  来源采集层：\texttt{source\_type}、\texttt{intel\_source}、\texttt{collection\_task}、\texttt{raw\_intel\_record}
\item
  攻击知识层：\texttt{attack\_entry}、\texttt{attack\_cvss\_assessment}、\texttt{attack\_evidence}、\texttt{attack\_taxonomy\_map}、\texttt{attack\_seed\_asset}、\texttt{remediation\_advice}
\item
  AI BOM
  层：\texttt{ai\_component}、\texttt{ai\_component\_alias}、\texttt{attack\_component\_impact}
\item
  治理审计层：\texttt{dedup\_audit}、\texttt{bom\_resolution\_queue}
\item
  读模型/视图：\texttt{v\_primary\_cvss\_score}、\texttt{v\_wp12\_attack\_feed}、\texttt{v\_component\_risk\_overview}、\texttt{v\_unresolved\_bom\_queue}、\texttt{v\_source\_quality\_dashboard}、\texttt{mv\_owasp\_coverage}
\end{itemize}

数据库模块需要解决的不是``怎么把所有 SQL 都塞进一个
repository.py''，而是：

\begin{enumerate}
\def\labelenumi{\arabic{enumi}.}
\tightlist
\item
  对采集 Agent、解析 Agent、BOM 解析 Agent、WP1-2 消费侧暴露清晰接口。
\item
  把事务边界、连接池、异常、模型映射从业务逻辑中抽离出来。
\item
  让 \texttt{agents/intel\_agents/}
  中的多个子智能体共用一套数据库访问规范。
\item
  保留 PostgreSQL
  特性（JSONB、GIN、视图、物化视图、\texttt{pg\_trgm}、\texttt{gen\_random\_uuid()}）的可用性，而不是被
  ORM 过度抽象掉。
\end{enumerate}

\begin{center}\rule{0.5\linewidth}{0.5pt}\end{center}

\hypertarget{ux8bbeux8ba1ux539fux5219}{%
\subsection{2. 设计原则}\label{ux8bbeux8ba1ux539fux5219}}

\hypertarget{ux91c7ux7528-psycopg3-dataclasspydantic-repository-unit-of-work}{%
\subsubsection{\texorpdfstring{2.1 采用
\texttt{psycopg3\ +\ dataclass/Pydantic\ +\ Repository\ +\ Unit\ of\ Work}}{2.1 采用 psycopg3 + dataclass/Pydantic + Repository + Unit of Work}}\label{ux91c7ux7528-psycopg3-dataclasspydantic-repository-unit-of-work}}

不建议把当前项目设计成``ORM-first''。原因如下：

\begin{itemize}
\tightlist
\item
  当前 schema 已经相当明确，且使用了 PostgreSQL 特性；
\item
  表、视图、物化视图较多，读写模式并不均匀；
\item
  很多对象天然更适合 \textbf{SQL 驱动}，例如：

  \begin{itemize}
  \tightlist
  \item
    基于视图的 WP1-2 投喂读取；
  \item
    去重审计与队列更新；
  \item
    BOM 模糊匹配；
  \item
    dashboard / overview 聚合读取。
  \end{itemize}
\end{itemize}

因此建议：

\begin{itemize}
\tightlist
\item
  \textbf{底层驱动}：\texttt{psycopg}（psycopg3）
\item
  \textbf{连接池}：\texttt{psycopg\_pool.ConnectionPool}
\item
  \textbf{领域模型}：\texttt{dataclass} 和
  \texttt{pydantic}（\texttt{pydantic}
  仅用于边界输入校验，内部数据都是用 \texttt{dataclass}）
\item
  \textbf{数据访问层}：分域 Repository
\item
  \textbf{事务管理}：Unit of Work
\item
  \textbf{服务层}：跨表写入编排，不直接暴露给 SQL 层
\end{itemize}

\hypertarget{repository-ux6309ux9886ux57dfux5b50ux57dfux62c6ux5206ux4e0dux6309-crud-ux751fux786cux62c6ux5206}{%
\subsubsection{2.2 Repository 按``领域子域''拆分，不按 CRUD
生硬拆分}\label{repository-ux6309ux9886ux57dfux5b50ux57dfux62c6ux5206ux4e0dux6309-crud-ux751fux786cux62c6ux5206}}

不要使用单一的 \texttt{repository.py} 承载所有数据库操作。当前 schema
至少应拆为：

\begin{itemize}
\tightlist
\item
  采集域 Repository
\item
  攻击知识域 Repository
\item
  AI BOM 域 Repository
\item
  治理审计域 Repository
\item
  读模型 / 视图 Repository
\end{itemize}

\hypertarget{ux8bfbux5199ux5206ux79bb}{%
\subsubsection{2.3 读写分离}\label{ux8bfbux5199ux5206ux79bb}}

当前数据库对象天然存在两类访问路径：

\begin{itemize}
\tightlist
\item
  \textbf{写路径}：面向表，强调事务一致性
\item
  \textbf{读路径}：面向视图 / 聚合视图，强调查询便利性
\end{itemize}

例如：

\begin{itemize}
\tightlist
\item
  \texttt{attack\_entry} / \texttt{attack\_cvss\_assessment} /
  \texttt{attack\_evidence} 是写路径
\item
  \texttt{v\_wp12\_attack\_feed} / \texttt{v\_component\_risk\_overview}
  是读路径
\end{itemize}

\hypertarget{ux6240ux6709ux8de8ux8868ux5199ux5165ux5fc5ux987bux7ecfux670dux52a1ux5c42ux6216-uow-ux5b8cux6210}{%
\subsubsection{2.4 所有跨表写入必须经服务层或 UoW
完成}\label{ux6240ux6709ux8de8ux8868ux5199ux5165ux5fc5ux987bux7ecfux670dux52a1ux5c42ux6216-uow-ux5b8cux6210}}

例如一次``原始情报解析入库''通常同时涉及：

\begin{itemize}
\tightlist
\item
  写 \texttt{raw\_intel\_record}
\item
  写 \texttt{attack\_entry}
\item
  写 \texttt{attack\_evidence}
\item
  写 \texttt{attack\_cvss\_assessment}
\item
  写 \texttt{attack\_taxonomy\_map}
\item
  写 \texttt{attack\_component\_impact}
\item
  必要时写 \texttt{bom\_resolution\_queue}
\end{itemize}

这种逻辑不能散落在 Agent 代码里，必须集中在 \textbf{service +
unit\_of\_work} 中。

\begin{center}\rule{0.5\linewidth}{0.5pt}\end{center}

\hypertarget{ux63a8ux8350ux76eeux5f55ux7ed3ux6784}{%
\subsection{3.
推荐目录结构}\label{ux63a8ux8350ux76eeux5f55ux7ed3ux6784}}

\begin{verbatim}
saads/
│
├─ agents/
│   └─ intel_agents/ （先不考虑这一部分）
│       ├─ collector_agent.py
│       ├─ parser_agent.py
│       ├─ cvss_agent.py
│       ├─ taxonomy_agent.py
│       ├─ bom_mapper_agent.py
│       └─ seed_builder_agent.py
│
├─ db/
│   ├─ __init__.py
│   ├─ connection.py
│   ├─ session.py
│   ├─ exceptions.py
│   ├─ unit_of_work.py
│   ├─ pagination.py
│   ├─ typing.py
│   ├─ sql/
│   │   ├─ attack_queries.py
│   │   ├─ bom_queries.py
│   │   ├─ source_queries.py
│   │   └─ governance_queries.py
│   ├─ models/
│   │   ├─ __init__.py
│   │   ├─ source.py
│   │   ├─ attack.py
│   │   ├─ component.py
│   │   ├─ governance.py
│   │   └─ views.py
│   ├─ repositories/
│   │   ├─ __init__.py
│   │   ├─ base.py
│   │   ├─ source_repository.py
│   │   ├─ attack_repository.py
│   │   ├─ component_repository.py
│   │   ├─ governance_repository.py
│   │   └─ read_model_repository.py
│   └─ services/
│       ├─ __init__.py
│       ├─ ingestion_service.py
│       ├─ attack_merge_service.py
│       ├─ cvss_service.py
│       ├─ taxonomy_service.py
│       ├─ bom_resolution_service.py
│       └─ wp12_feed_service.py
│
└─ config.py
\end{verbatim}

这个结构里，最关键的是：

\begin{itemize}
\tightlist
\item
  \texttt{models/}：定义 Python 侧数据对象
\item
  \texttt{repositories/}：封装单域数据库访问
\item
  \texttt{services/}：封装跨表业务事务
\item
  \texttt{unit\_of\_work.py}：统一事务边界
\item
  \texttt{read\_model\_repository.py}：专门读视图和物化视图
\end{itemize}

\begin{center}\rule{0.5\linewidth}{0.5pt}\end{center}

\hypertarget{config.py-ux4e0eux6570ux636eux5e93ux914dux7f6e}{%
\subsection{\texorpdfstring{4. \texttt{config.py}
与数据库配置}{4. config.py 与数据库配置}}\label{config.py-ux4e0eux6570ux636eux5e93ux914dux7f6e}}

数据库模块不应自己硬编码连接参数。建议把配置统一放在项目根部
\texttt{config.py}。

建议配置项：

\begin{Shaded}
\begin{Highlighting}[]
\NormalTok{POSTGRES\_HOST}
\NormalTok{POSTGRES\_PORT}
\NormalTok{POSTGRES\_DB}
\NormalTok{POSTGRES\_USER}
\NormalTok{POSTGRES\_PASSWORD}
\NormalTok{POSTGRES\_SCHEMA }\OperatorTok{=} \StringTok{"wp11"}
\NormalTok{POSTGRES\_MIN\_SIZE }\OperatorTok{=} \DecValTok{1}
\NormalTok{POSTGRES\_MAX\_SIZE }\OperatorTok{=} \DecValTok{10}
\NormalTok{POSTGRES\_CONNECT\_TIMEOUT }\OperatorTok{=} \DecValTok{5}
\NormalTok{POSTGRES\_STATEMENT\_TIMEOUT\_MS }\OperatorTok{=} \DecValTok{30000}
\end{Highlighting}
\end{Shaded}

也可以补充：

\begin{Shaded}
\begin{Highlighting}[]
\NormalTok{POSTGRES\_DSN}
\end{Highlighting}
\end{Shaded}

推荐优先使用 DSN。

\begin{center}\rule{0.5\linewidth}{0.5pt}\end{center}

\hypertarget{dbconnection.pyux8fdeux63a5ux6c60ux4e0eux8fdeux63a5ux5de5ux5382}{%
\subsection{\texorpdfstring{5.
\texttt{db/connection.py}：连接池与连接工厂}{5. db/connection.py：连接池与连接工厂}}\label{dbconnection.pyux8fdeux63a5ux6c60ux4e0eux8fdeux63a5ux5de5ux5382}}

\hypertarget{ux804cux8d23}{%
\subsubsection{职责}\label{ux804cux8d23}}

\begin{itemize}
\tightlist
\item
  创建和持有 PostgreSQL 连接池
\item
  为上层提供连接获取接口
\item
  在连接建立后设置：

  \begin{itemize}
  \tightlist
  \item
    \texttt{search\_path\ =\ wp11,\ public}
  \item
    statement timeout
  \item
    row factory（dict row 或 class row）
  \end{itemize}
\end{itemize}

\hypertarget{ux8bbeux8ba1ux8981ux70b9}{%
\subsubsection{设计要点}\label{ux8bbeux8ba1ux8981ux70b9}}

\begin{enumerate}
\def\labelenumi{\arabic{enumi}.}
\tightlist
\item
  使用全局单例连接池，但初始化动作应显式完成。
\item
  避免在 import 时自动建连。
\item
  提供同步上下文管理接口。
\end{enumerate}

\hypertarget{ux63a8ux8350ux63a5ux53e3}{%
\subsubsection{推荐接口}\label{ux63a8ux8350ux63a5ux53e3}}

\begin{itemize}
\tightlist
\item
  \texttt{init\_pool()}
\item
  \texttt{close\_pool()}
\item
  \texttt{get\_connection()}
\item
  \texttt{connection\_context()}
\end{itemize}

\hypertarget{ux4e0dux8981ux505aux7684ux4e8b}{%
\subsubsection{不要做的事}\label{ux4e0dux8981ux505aux7684ux4e8b}}

\begin{itemize}
\tightlist
\item
  不要在每个 repository 里自己 \texttt{psycopg.connect(...)}
\item
  不要在 Agent 里直接管理连接生命周期
\end{itemize}

\begin{center}\rule{0.5\linewidth}{0.5pt}\end{center}

\hypertarget{dbsession.pyux6e38ux6807ux4e0eux6267ux884cux8f85ux52a9}{%
\subsection{\texorpdfstring{6.
\texttt{db/session.py}：游标与执行辅助}{6. db/session.py：游标与执行辅助}}\label{dbsession.pyux6e38ux6807ux4e0eux6267ux884cux8f85ux52a9}}

\hypertarget{ux804cux8d23-1}{%
\subsubsection{职责}\label{ux804cux8d23-1}}

\begin{itemize}
\tightlist
\item
  统一 cursor 的创建方式
\item
  统一 SQL 执行的参数绑定与日志记录
\item
  提供常见执行封装：

  \begin{itemize}
  \tightlist
  \item
    \texttt{fetch\_one}
  \item
    \texttt{fetch\_all}
  \item
    \texttt{execute}
  \item
    \texttt{execute\_many}
  \item
    \texttt{fetch\_scalar}
  \end{itemize}
\end{itemize}

\hypertarget{ux4e3aux4ec0ux4e48ux9700ux8981ux8fd9ux4e00ux5c42}{%
\subsubsection{为什么需要这一层}\label{ux4e3aux4ec0ux4e48ux9700ux8981ux8fd9ux4e00ux5c42}}

因为当前系统里会有大量``读一个视图''``按条件更新一个队列状态''``执行
upsert''的小型操作。如果所有 repository 都自己处理
cursor，会重复很多错误处理和日志代码。

\begin{center}\rule{0.5\linewidth}{0.5pt}\end{center}

\hypertarget{dbexceptions.pyux5f02ux5e38ux4f53ux7cfb}{%
\subsection{\texorpdfstring{7.
\texttt{db/exceptions.py}：异常体系}{7. db/exceptions.py：异常体系}}\label{dbexceptions.pyux5f02ux5e38ux4f53ux7cfb}}

建议定义数据库模块自己的异常层级，避免上层 Agent 直接依赖 psycopg
原始异常。

\hypertarget{ux63a8ux8350ux5f02ux5e38ux7c7bux578b}{%
\subsubsection{推荐异常类型}\label{ux63a8ux8350ux5f02ux5e38ux7c7bux578b}}

\begin{itemize}
\tightlist
\item
  \texttt{DatabaseError}：数据库模块基类异常
\item
  \texttt{ConnectionInitError}
\item
  \texttt{QueryExecutionError}
\item
  \texttt{NotFoundError}
\item
  \texttt{DuplicateEntityError}
\item
  \texttt{ValidationError}
\item
  \texttt{ConcurrencyConflictError}
\item
  \texttt{TransactionError}
\end{itemize}

\hypertarget{ux8bbeux8ba1ux539fux5219-1}{%
\subsubsection{设计原则}\label{ux8bbeux8ba1ux539fux5219-1}}

\begin{itemize}
\tightlist
\item
  psycopg 的异常在 repository 内部转换为项目级异常
\item
  Agent 只处理项目级异常，不感知底层驱动细节
\end{itemize}

\begin{center}\rule{0.5\linewidth}{0.5pt}\end{center}

\hypertarget{dbunit_of_work.pyux4e8bux52a1ux8fb9ux754cux6838ux5fc3}{%
\subsection{\texorpdfstring{8.
\texttt{db/unit\_of\_work.py}：事务边界核心}{8. db/unit\_of\_work.py：事务边界核心}}\label{dbunit_of_work.pyux4e8bux52a1ux8fb9ux754cux6838ux5fc3}}

\hypertarget{ux4f5cux7528}{%
\subsubsection{作用}\label{ux4f5cux7528}}

统一管理跨 Repository 的事务。

\hypertarget{ux4e3aux4ec0ux4e48ux5fc5ux987bux6709-uow}{%
\subsubsection{为什么必须有
UoW}\label{ux4e3aux4ec0ux4e48ux5fc5ux987bux6709-uow}}

下面这些动作天然是``一次事务''：

\begin{enumerate}
\def\labelenumi{\arabic{enumi}.}
\tightlist
\item
  新建 \texttt{attack\_entry}
\item
  写入证据 \texttt{attack\_evidence}
\item
  写入主 CVSS \texttt{attack\_cvss\_assessment}
\item
  写入 taxonomy \texttt{attack\_taxonomy\_map}
\item
  写入组件影响 \texttt{attack\_component\_impact}
\item
  必要时创建 \texttt{bom\_resolution\_queue}
\end{enumerate}

如果中间任何一步失败，前面的写入必须回滚。

\hypertarget{ux63a8ux8350ux80fdux529b}{%
\subsubsection{推荐能力}\label{ux63a8ux8350ux80fdux529b}}

\begin{itemize}
\tightlist
\item
  \texttt{\_\_enter\_\_\ /\ \_\_exit\_\_}
\item
  自动 commit / rollback
\item
  暴露按域组织的 repositories，例如：

  \begin{itemize}
  \tightlist
  \item
    \texttt{uow.sources}
  \item
    \texttt{uow.attacks}
  \item
    \texttt{uow.components}
  \item
    \texttt{uow.governance}
  \item
    \texttt{uow.read\_models}
  \end{itemize}
\end{itemize}

\hypertarget{ux4e00ux4e2a-uow-ux4e2dux4e0dux5e94ux505aux7684ux4e8b}{%
\subsubsection{一个 UoW
中不应做的事}\label{ux4e00ux4e2a-uow-ux4e2dux4e0dux5e94ux505aux7684ux4e8b}}

\begin{itemize}
\tightlist
\item
  不要在长时间网络请求外持有事务
\item
  采集 HTTP 拉取、LLM
  推理、规则计算等慢操作，应在事务外完成；事务内只做确定的数据库读写
\end{itemize}

\begin{center}\rule{0.5\linewidth}{0.5pt}\end{center}

\hypertarget{dbmodelspython-ux6570ux636eux6a21ux578bux8bbeux8ba1}{%
\subsection{\texorpdfstring{9. \texttt{db/models/}：Python
数据模型设计}{9. db/models/：Python 数据模型设计}}\label{dbmodelspython-ux6570ux636eux6a21ux578bux8bbeux8ba1}}

这一层负责从 PostgreSQL 的物理表映射到 Python 可操作对象。

推荐使用 \texttt{dataclass}
表达``数据库记录对象''。如果边界输入复杂，可额外定义 \texttt{pydantic}
输入 DTO。

\begin{center}\rule{0.5\linewidth}{0.5pt}\end{center}

\hypertarget{dbmodelssource.py}{%
\subsection{\texorpdfstring{10.
\texttt{db/models/source.py}}{10. db/models/source.py}}\label{dbmodelssource.py}}

\hypertarget{ux5efaux8baeux6a21ux578b}{%
\subsubsection{建议模型}\label{ux5efaux8baeux6a21ux578b}}

\begin{itemize}
\tightlist
\item
  \texttt{SourceType}
\item
  \texttt{IntelSource}
\item
  \texttt{CollectionTask}
\item
  \texttt{RawIntelRecord}
\end{itemize}

\hypertarget{ux4e3bux8981ux7528ux9014}{%
\subsubsection{主要用途}\label{ux4e3bux8981ux7528ux9014}}

\begin{itemize}
\tightlist
\item
  采集源注册
\item
  采集任务生命周期管理
\item
  原始情报入库
\end{itemize}

\hypertarget{ux7279ux522bux8bf4ux660e}{%
\subsubsection{特别说明}\label{ux7279ux522bux8bf4ux660e}}

\texttt{RawIntelRecord} 建议包含：

\begin{itemize}
\tightlist
\item
  \texttt{raw\_id}
\item
  \texttt{source\_id}
\item
  \texttt{task\_id}
\item
  \texttt{source\_uri}
\item
  \texttt{title}
\item
  \texttt{content\_hash}
\item
  \texttt{raw\_format}
\item
  \texttt{payload\_uri}
\item
  \texttt{language\_code}
\item
  \texttt{relevance\_score}
\item
  \texttt{parser\_status}
\item
  \texttt{fetched\_at}
\item
  \texttt{created\_at}
\item
  \texttt{is\_deleted}
\end{itemize}

\begin{center}\rule{0.5\linewidth}{0.5pt}\end{center}

\hypertarget{dbmodelsattack.py}{%
\subsection{\texorpdfstring{11.
\texttt{db/models/attack.py}}{11. db/models/attack.py}}\label{dbmodelsattack.py}}

\hypertarget{ux5efaux8baeux6a21ux578b-1}{%
\subsubsection{建议模型}\label{ux5efaux8baeux6a21ux578b-1}}

\begin{itemize}
\tightlist
\item
  \texttt{AttackEntry}
\item
  \texttt{AttackCvssAssessment}
\item
  \texttt{AttackEvidence}
\item
  \texttt{AttackTaxonomyMap}
\item
  \texttt{AttackSeedAsset}
\item
  \texttt{RemediationAdvice}
\end{itemize}

\hypertarget{ux8bbeux8ba1ux8bf4ux660e}{%
\subsubsection{设计说明}\label{ux8bbeux8ba1ux8bf4ux660e}}

\hypertarget{attackentry}{%
\paragraph{\texorpdfstring{\texttt{AttackEntry}}{AttackEntry}}\label{attackentry}}

对应标准化攻击知识主表，是整个 WP1-1 的核心模型。

\hypertarget{attackcvssassessment}{%
\paragraph{\texorpdfstring{\texttt{AttackCvssAssessment}}{AttackCvssAssessment}}\label{attackcvssassessment}}

必须独立为单独模型，不要并入 \texttt{AttackEntry}。原因：

\begin{itemize}
\tightlist
\item
  一条攻击条目可以有多条 CVSS 记录
\item
  同一攻击可能同时存在 \texttt{supplied} / \texttt{calculated} /
  \texttt{estimated} / \texttt{manual}
\item
  \texttt{is\_primary} 只标识对外主分数，不等于唯一分数
\end{itemize}

\hypertarget{attackevidence}{%
\paragraph{\texorpdfstring{\texttt{AttackEvidence}}{AttackEvidence}}\label{attackevidence}}

这是攻击条目与原始情报之间的证据链表，不能省略，否则
\texttt{attack\_entry} 会失去可追溯性。

\hypertarget{attackseedasset}{%
\paragraph{\texorpdfstring{\texttt{AttackSeedAsset}}{AttackSeedAsset}}\label{attackseedasset}}

这是 WP1-2 的关键供给对象，代表 PoC、payload 模板、prompt 语料、rule
等``测试种子''。

\begin{center}\rule{0.5\linewidth}{0.5pt}\end{center}

\hypertarget{dbmodelscomponent.py}{%
\subsection{\texorpdfstring{12.
\texttt{db/models/component.py}}{12. db/models/component.py}}\label{dbmodelscomponent.py}}

\hypertarget{ux5efaux8baeux6a21ux578b-2}{%
\subsubsection{建议模型}\label{ux5efaux8baeux6a21ux578b-2}}

\begin{itemize}
\tightlist
\item
  \texttt{AiComponent}
\item
  \texttt{AiComponentAlias}
\item
  \texttt{AttackComponentImpact}
\end{itemize}

\hypertarget{ux8bbeux8ba1ux8bf4ux660e-1}{%
\subsubsection{设计说明}\label{ux8bbeux8ba1ux8bf4ux660e-1}}

这一组模型服务 AI BOM 适配逻辑。

\hypertarget{aicomponent}{%
\paragraph{\texorpdfstring{\texttt{AiComponent}}{AiComponent}}\label{aicomponent}}

标准 AI 组件主表。

\hypertarget{aicomponentalias}{%
\paragraph{\texorpdfstring{\texttt{AiComponentAlias}}{AiComponentAlias}}\label{aicomponentalias}}

负责别名规范化与模糊匹配。由于 schema 上已经有
\texttt{normalized\_alias} + trigram 索引，这意味着 Python
侧必须明确区分：

\begin{itemize}
\tightlist
\item
  原始别名
\item
  规范化别名
\end{itemize}

\hypertarget{attackcomponentimpact}{%
\paragraph{\texorpdfstring{\texttt{AttackComponentImpact}}{AttackComponentImpact}}\label{attackcomponentimpact}}

用于记录``某攻击影响哪些组件以及以何种版本约束影响''。这是 AI BOM
风险联动的关键桥表。

\begin{center}\rule{0.5\linewidth}{0.5pt}\end{center}

\hypertarget{dbmodelsgovernance.py}{%
\subsection{\texorpdfstring{13.
\texttt{db/models/governance.py}}{13. db/models/governance.py}}\label{dbmodelsgovernance.py}}

\hypertarget{ux5efaux8baeux6a21ux578b-3}{%
\subsubsection{建议模型}\label{ux5efaux8baeux6a21ux578b-3}}

\begin{itemize}
\tightlist
\item
  \texttt{DedupAudit}
\item
  \texttt{BomResolutionQueueItem}
\end{itemize}

\hypertarget{ux4f5cux7528-1}{%
\subsubsection{作用}\label{ux4f5cux7528-1}}

治理审计模型不是附属信息，而是运营层对象：

\begin{itemize}
\tightlist
\item
  \texttt{DedupAudit}：记录去重决策
\item
  \texttt{BomResolutionQueueItem}：记录 AI BOM
  未解析项的人工/半自动复核队列
\end{itemize}

这些对象通常会被独立的治理 Agent 或后台管理界面消费。

\begin{center}\rule{0.5\linewidth}{0.5pt}\end{center}

\hypertarget{dbmodelsviews.py}{%
\subsection{\texorpdfstring{14.
\texttt{db/models/views.py}}{14. db/models/views.py}}\label{dbmodelsviews.py}}

建议把视图映射对象集中在这里：

\begin{itemize}
\tightlist
\item
  \texttt{PrimaryCvssView}
\item
  \texttt{Wp12AttackFeedRow}
\item
  \texttt{ComponentRiskOverviewRow}
\item
  \texttt{UnresolvedBomQueueRow}
\item
  \texttt{SourceQualityDashboardRow}
\item
  \texttt{OwaspCoverageRow}
\end{itemize}

\hypertarget{ux8bbeux8ba1ux539fux5219-2}{%
\subsubsection{设计原则}\label{ux8bbeux8ba1ux539fux5219-2}}

\begin{itemize}
\tightlist
\item
  视图模型是只读对象
\item
  不参与写回
\item
  单独放置，避免与表模型混淆
\end{itemize}

\begin{center}\rule{0.5\linewidth}{0.5pt}\end{center}

\hypertarget{dbrepositoriesbase.py}{%
\subsection{\texorpdfstring{15.
\texttt{db/repositories/base.py}}{15. db/repositories/base.py}}\label{dbrepositoriesbase.py}}

\hypertarget{ux4f5cux7528-2}{%
\subsubsection{作用}\label{ux4f5cux7528-2}}

封装所有 repository 的公共能力：

\begin{itemize}
\tightlist
\item
  持有 connection / session
\item
  执行 SQL
\item
  统一日志与异常转换
\item
  常见辅助函数（唯一查找、分页、批量插入）
\end{itemize}

\hypertarget{ux5efaux8baeux516cux5171ux65b9ux6cd5}{%
\subsubsection{建议公共方法}\label{ux5efaux8baeux516cux5171ux65b9ux6cd5}}

\begin{itemize}
\tightlist
\item
  \texttt{\_fetch\_one(...)}
\item
  \texttt{\_fetch\_all(...)}
\item
  \texttt{\_execute(...)}
\item
  \texttt{\_execute\_many(...)}
\item
  \texttt{\_fetch\_scalar(...)}
\end{itemize}

\begin{center}\rule{0.5\linewidth}{0.5pt}\end{center}

\hypertarget{source_repository.py}{%
\subsection{\texorpdfstring{16.
\texttt{source\_repository.py}}{16. source\_repository.py}}\label{source_repository.py}}

\hypertarget{ux8d1fux8d23ux5bf9ux8c61}{%
\subsubsection{负责对象}\label{ux8d1fux8d23ux5bf9ux8c61}}

\begin{itemize}
\tightlist
\item
  \texttt{source\_type}
\item
  \texttt{intel\_source}
\item
  \texttt{collection\_task}
\item
  \texttt{raw\_intel\_record}
\end{itemize}

\hypertarget{ux5efaux8baeux63a5ux53e3}{%
\subsubsection{建议接口}\label{ux5efaux8baeux63a5ux53e3}}

\begin{itemize}
\tightlist
\item
  \texttt{get\_source\_by\_name(source\_name)}
\item
  \texttt{list\_enabled\_sources(source\_type=None)}
\item
  \texttt{create\_collection\_task(...)}
\item
  \texttt{update\_collection\_task\_status(...)}
\item
  \texttt{insert\_raw\_intel\_record(...)}
\item
  \texttt{get\_raw\_record\_by\_hash(source\_id,\ content\_hash)}
\item
  \texttt{mark\_raw\_record\_parser\_status(raw\_id,\ status)}
\item
  \texttt{list\_pending\_raw\_records(limit)}
\end{itemize}

\hypertarget{ux4e1aux52a1ux7279ux70b9}{%
\subsubsection{业务特点}\label{ux4e1aux52a1ux7279ux70b9}}

这一层是采集智能体与数据库的第一接触面。重点是：

\begin{itemize}
\tightlist
\item
  控制重复原始数据写入
\item
  管理 task 生命周期
\item
  为解析 Agent 提供待处理原始记录
\end{itemize}

\begin{center}\rule{0.5\linewidth}{0.5pt}\end{center}

\hypertarget{attack_repository.py}{%
\subsection{\texorpdfstring{17.
\texttt{attack\_repository.py}}{17. attack\_repository.py}}\label{attack_repository.py}}

\hypertarget{ux8d1fux8d23ux5bf9ux8c61-1}{%
\subsubsection{负责对象}\label{ux8d1fux8d23ux5bf9ux8c61-1}}

\begin{itemize}
\tightlist
\item
  \texttt{attack\_entry}
\item
  \texttt{attack\_cvss\_assessment}
\item
  \texttt{attack\_evidence}
\item
  \texttt{attack\_taxonomy\_map}
\item
  \texttt{attack\_seed\_asset}
\item
  \texttt{remediation\_advice}
\end{itemize}

\hypertarget{ux5efaux8baeux63a5ux53e3-1}{%
\subsubsection{建议接口}\label{ux5efaux8baeux63a5ux53e3-1}}

\begin{itemize}
\tightlist
\item
  \texttt{get\_attack\_by\_code(attack\_code)}
\item
  \texttt{create\_attack\_entry(...)}
\item
  \texttt{update\_attack\_entry(...)}
\item
  \texttt{upsert\_attack\_entry\_by\_code(...)}
\item
  \texttt{insert\_attack\_evidence(...)}
\item
  \texttt{list\_attack\_evidence(attack\_id)}
\item
  \texttt{insert\_cvss\_assessment(...)}
\item
  \texttt{list\_cvss\_assessments(attack\_id)}
\item
  \texttt{set\_primary\_cvss(score\_id)}
\item
  \texttt{replace\_primary\_taxonomy(...)}
\item
  \texttt{insert\_seed\_asset(...)}
\item
  \texttt{list\_published\_seed\_assets(attack\_id)}
\item
  \texttt{insert\_remediation\_advice(...)}
\end{itemize}

\hypertarget{ux8bbeux8ba1ux7ec6ux8282}{%
\subsubsection{设计细节}\label{ux8bbeux8ba1ux7ec6ux8282}}

\hypertarget{ux5173ux4e8e-set_primary_cvssscore_id}{%
\paragraph{\texorpdfstring{关于
\texttt{set\_primary\_cvss(score\_id)}}{关于 set\_primary\_cvss(score\_id)}}\label{ux5173ux4e8e-set_primary_cvssscore_id}}

由于 schema 上存在
\texttt{uq\_cvss\_primary\_per\_attack\_version}，因此切换主 CVSS
时必须：

\begin{enumerate}
\def\labelenumi{\arabic{enumi}.}
\tightlist
\item
  先找到目标记录的 \texttt{attack\_id} 与 \texttt{cvss\_version}
\item
  把同攻击、同版本的其他记录置为 \texttt{is\_primary\ =\ false}
\item
  再把目标记录置为 \texttt{true}
\end{enumerate}

必须在同一事务中执行。

\hypertarget{ux5173ux4e8e-taxonomy}{%
\paragraph{关于 taxonomy}\label{ux5173ux4e8e-taxonomy}}

\texttt{attack\_taxonomy\_map} 同样有``每种 taxonomy\_type
只能一个主项''的约束，因此 \texttt{replace\_primary\_taxonomy(...)}
也应走事务。

\begin{center}\rule{0.5\linewidth}{0.5pt}\end{center}

\hypertarget{component_repository.py}{%
\subsection{\texorpdfstring{18.
\texttt{component\_repository.py}}{18. component\_repository.py}}\label{component_repository.py}}

\hypertarget{ux8d1fux8d23ux5bf9ux8c61-2}{%
\subsubsection{负责对象}\label{ux8d1fux8d23ux5bf9ux8c61-2}}

\begin{itemize}
\tightlist
\item
  \texttt{ai\_component}
\item
  \texttt{ai\_component\_alias}
\item
  \texttt{attack\_component\_impact}
\end{itemize}

\hypertarget{ux5efaux8baeux63a5ux53e3-2}{%
\subsubsection{建议接口}\label{ux5efaux8baeux63a5ux53e3-2}}

\begin{itemize}
\tightlist
\item
  \texttt{get\_component\_by\_code(component\_code)}
\item
  \texttt{get\_component\_by\_name(component\_name)}
\item
  \texttt{search\_component\_alias(normalized\_alias,\ limit=10)}
\item
  \texttt{create\_component(...)}
\item
  \texttt{insert\_component\_alias(...)}
\item
  \texttt{list\_component\_aliases(component\_id)}
\item
  \texttt{upsert\_attack\_component\_impact(...)}
\item
  \texttt{list\_component\_impacts\_by\_attack(attack\_id)}
\item
  \texttt{list\_attacks\_by\_component(component\_id)}
\end{itemize}

\hypertarget{ux8bbeux8ba1ux7ec6ux8282-1}{%
\subsubsection{设计细节}\label{ux8bbeux8ba1ux7ec6ux8282-1}}

由于 \texttt{ai\_component\_alias.normalized\_alias}
是全局唯一，别名标准化函数必须稳定，否则会导致重复别名和解析错误。

建议在 Python 层统一实现：

\begin{itemize}
\tightlist
\item
  小写化
\item
  去空格 / 下划线 / 连字符标准化
\item
  去 vendor 前后冗余前缀
\end{itemize}

\begin{center}\rule{0.5\linewidth}{0.5pt}\end{center}

\hypertarget{governance_repository.py}{%
\subsection{\texorpdfstring{19.
\texttt{governance\_repository.py}}{19. governance\_repository.py}}\label{governance_repository.py}}

\hypertarget{ux8d1fux8d23ux5bf9ux8c61-3}{%
\subsubsection{负责对象}\label{ux8d1fux8d23ux5bf9ux8c61-3}}

\begin{itemize}
\tightlist
\item
  \texttt{dedup\_audit}
\item
  \texttt{bom\_resolution\_queue}
\end{itemize}

\hypertarget{ux5efaux8baeux63a5ux53e3-3}{%
\subsubsection{建议接口}\label{ux5efaux8baeux63a5ux53e3-3}}

\begin{itemize}
\tightlist
\item
  \texttt{insert\_dedup\_audit(...)}
\item
  \texttt{list\_dedup\_candidates\_for\_review(...)}
\item
  \texttt{enqueue\_bom\_resolution(...)}
\item
  \texttt{list\_open\_bom\_queue(limit)}
\item
  \texttt{resolve\_bom\_queue\_item(queue\_id,\ resolved\_component\_id)}
\item
  \texttt{reject\_bom\_queue\_item(queue\_id)}
\end{itemize}

\hypertarget{ux8bbeux8ba1ux8bf4ux660e-2}{%
\subsubsection{设计说明}\label{ux8bbeux8ba1ux8bf4ux660e-2}}

这里不应只做``后台管理接口''。对于 \texttt{bom\_mapper\_agent}
来说，无法自动解析的组件必须立即写入 queue，而不是静默失败。

\begin{center}\rule{0.5\linewidth}{0.5pt}\end{center}

\hypertarget{read_model_repository.py}{%
\subsection{\texorpdfstring{20.
\texttt{read\_model\_repository.py}}{20. read\_model\_repository.py}}\label{read_model_repository.py}}

\hypertarget{ux8d1fux8d23ux5bf9ux8c61-4}{%
\subsubsection{负责对象}\label{ux8d1fux8d23ux5bf9ux8c61-4}}

\begin{itemize}
\tightlist
\item
  \texttt{v\_primary\_cvss\_score}
\item
  \texttt{v\_wp12\_attack\_feed}
\item
  \texttt{v\_component\_risk\_overview}
\item
  \texttt{v\_unresolved\_bom\_queue}
\item
  \texttt{v\_source\_quality\_dashboard}
\item
  \texttt{mv\_owasp\_coverage}
\end{itemize}

\hypertarget{ux5efaux8baeux63a5ux53e3-4}{%
\subsubsection{建议接口}\label{ux5efaux8baeux63a5ux53e3-4}}

\begin{itemize}
\tightlist
\item
  \texttt{get\_primary\_cvss(attack\_id)}
\item
  \texttt{list\_wp12\_attack\_feed(filters...)}
\item
  \texttt{list\_component\_risk\_overview(...)}
\item
  \texttt{list\_unresolved\_bom\_queue(...)}
\item
  \texttt{get\_source\_quality\_dashboard(...)}
\item
  \texttt{list\_owasp\_coverage(...)}
\item
  \texttt{refresh\_mv\_owasp\_coverage()}
\end{itemize}

\hypertarget{ux4e3aux4ec0ux4e48ux5355ux72ecux62c6ux5206}{%
\subsubsection{为什么单独拆分}\label{ux4e3aux4ec0ux4e48ux5355ux72ecux62c6ux5206}}

因为这些查询面向消费方：

\begin{itemize}
\tightlist
\item
  WP1-2
\item
  dashboard
\item
  风险总览
\item
  运营分析
\end{itemize}

它们的目标不是保持表写入一致性，而是提供稳定的读模型。

\begin{center}\rule{0.5\linewidth}{0.5pt}\end{center}

\hypertarget{dbservicesux670dux52a1ux5c42ux8bbeux8ba1}{%
\subsection{\texorpdfstring{21.
\texttt{db/services/}：服务层设计}{21. db/services/：服务层设计}}\label{dbservicesux670dux52a1ux5c42ux8bbeux8ba1}}

服务层负责``跨表业务动作''的编排。建议至少有以下服务：

\hypertarget{ingestion_service.py}{%
\subsubsection{\texorpdfstring{21.1
\texttt{ingestion\_service.py}}{21.1 ingestion\_service.py}}\label{ingestion_service.py}}

负责：

\begin{itemize}
\tightlist
\item
  创建采集任务
\item
  原始记录入库
\item
  防重复检查
\item
  parser 状态更新
\end{itemize}

\hypertarget{attack_merge_service.py}{%
\subsubsection{\texorpdfstring{21.2
\texttt{attack\_merge\_service.py}}{21.2 attack\_merge\_service.py}}\label{attack_merge_service.py}}

负责：

\begin{itemize}
\tightlist
\item
  将解析结果合并成 \texttt{attack\_entry}
\item
  写证据链
\item
  建立去重审计记录
\end{itemize}

\hypertarget{cvss_service.py}{%
\subsubsection{\texorpdfstring{21.3
\texttt{cvss\_service.py}}{21.3 cvss\_service.py}}\label{cvss_service.py}}

负责：

\begin{itemize}
\tightlist
\item
  写入 CVSS
\item
  计算/补全分数
\item
  切换主评分
\end{itemize}

\hypertarget{taxonomy_service.py}{%
\subsubsection{\texorpdfstring{21.4
\texttt{taxonomy\_service.py}}{21.4 taxonomy\_service.py}}\label{taxonomy_service.py}}

负责：

\begin{itemize}
\tightlist
\item
  批量写 taxonomy
\item
  替换主 taxonomy
\end{itemize}

\hypertarget{bom_resolution_service.py}{%
\subsubsection{\texorpdfstring{21.5
\texttt{bom\_resolution\_service.py}}{21.5 bom\_resolution\_service.py}}\label{bom_resolution_service.py}}

负责：

\begin{itemize}
\tightlist
\item
  组件别名匹配
\item
  自动解析成功则写 \texttt{attack\_component\_impact}
\item
  自动解析失败则写 \texttt{bom\_resolution\_queue}
\end{itemize}

\hypertarget{wp12_feed_service.py}{%
\subsubsection{\texorpdfstring{21.6
\texttt{wp12\_feed\_service.py}}{21.6 wp12\_feed\_service.py}}\label{wp12_feed_service.py}}

负责：

\begin{itemize}
\tightlist
\item
  面向 WP1-2 构造``测试脚本生成输入集''
\item
  优先通过 \texttt{v\_wp12\_attack\_feed} 提供读取接口
\end{itemize}

\begin{center}\rule{0.5\linewidth}{0.5pt}\end{center}

\hypertarget{sql-ux7ec4ux7ec7ux65b9ux5f0fdbsql}{%
\subsection{\texorpdfstring{22. SQL
组织方式：\texttt{db/sql/}}{22. SQL 组织方式：db/sql/}}\label{sql-ux7ec4ux7ec7ux65b9ux5f0fdbsql}}

不建议把所有 SQL 内联进 repository 方法体。建议把复杂 SQL 放入
\texttt{db/sql/*.py} 中按领域维护。

\hypertarget{ux4f18ux70b9}{%
\subsubsection{优点}\label{ux4f18ux70b9}}

\begin{itemize}
\tightlist
\item
  SQL 可测试
\item
  SQL 与控制流解耦
\item
  复杂查询（尤其视图查询、模糊匹配、聚合）更易维护
\end{itemize}

\hypertarget{ux9002ux5408ux62bdux79bbux7684-sql}{%
\subsubsection{适合抽离的
SQL}\label{ux9002ux5408ux62bdux79bbux7684-sql}}

\begin{itemize}
\tightlist
\item
  feed 查询
\item
  alias 模糊匹配
\item
  dedup 候选查询
\item
  dashboard 聚合读取
\end{itemize}

\begin{center}\rule{0.5\linewidth}{0.5pt}\end{center}

\hypertarget{ux4e8bux52a1ux8fb9ux754cux5efaux8bae}{%
\subsection{23.
事务边界建议}\label{ux4e8bux52a1ux8fb9ux754cux5efaux8bae}}

\hypertarget{ux4e00ux6761ux539fux59cbux60c5ux62a5ux6293ux53d6ux4efbux52a1}{%
\subsubsection{23.1
一条原始情报抓取任务}\label{ux4e00ux6761ux539fux59cbux60c5ux62a5ux6293ux53d6ux4efbux52a1}}

应拆成两个阶段：

\hypertarget{ux9636ux6bb5-aux4e8bux52a1ux5916}{%
\paragraph{阶段 A：事务外}\label{ux9636ux6bb5-aux4e8bux52a1ux5916}}

\begin{itemize}
\tightlist
\item
  发 HTTP 请求
\item
  下载 payload
\item
  算 hash
\item
  解析内容
\item
  调用模型 / LLM
\end{itemize}

\hypertarget{ux9636ux6bb5-bux4e8bux52a1ux5185}{%
\paragraph{阶段 B：事务内}\label{ux9636ux6bb5-bux4e8bux52a1ux5185}}

\begin{itemize}
\tightlist
\item
  插入 \texttt{collection\_task} 更新状态
\item
  插入 \texttt{raw\_intel\_record}
\item
  插入 / 更新 \texttt{attack\_entry}
\item
  插入 \texttt{attack\_evidence}
\item
  插入 \texttt{attack\_cvss\_assessment}
\item
  插入 taxonomy / component impact
\item
  更新 task 状态
\end{itemize}

\hypertarget{ux4e0dux5efaux8baeux5728ux4e00ux4e2aux4e8bux52a1ux4e2dux540cux65f6ux505a}{%
\subsubsection{23.2
不建议在一个事务中同时做}\label{ux4e0dux5efaux8baeux5728ux4e00ux4e2aux4e8bux52a1ux4e2dux540cux65f6ux505a}}

\begin{itemize}
\tightlist
\item
  网络访问
\item
  长时间文件 IO
\item
  LLM 推理
\item
  数据库写入
\end{itemize}

否则锁持有时间过长。

\begin{center}\rule{0.5\linewidth}{0.5pt}\end{center}

\hypertarget{ux5e76ux53d1ux4e0eux5e42ux7b49ux6027ux8bbeux8ba1}{%
\subsection{24.
并发与幂等性设计}\label{ux5e76ux53d1ux4e0eux5e42ux7b49ux6027ux8bbeux8ba1}}

\hypertarget{ux539fux59cbux8bb0ux5f55ux5e42ux7b49}{%
\subsubsection{24.1
原始记录幂等}\label{ux539fux59cbux8bb0ux5f55ux5e42ux7b49}}

基于：

\begin{itemize}
\tightlist
\item
  \texttt{raw\_intel\_record\ (source\_id,\ content\_hash)} 唯一约束
\end{itemize}

因此 Repository 应优先提供：

\begin{itemize}
\tightlist
\item
  \texttt{insert\_or\_get\_raw\_record(...)}
\end{itemize}

\hypertarget{ux522bux540dux5e42ux7b49}{%
\subsubsection{24.2 别名幂等}\label{ux522bux540dux5e42ux7b49}}

基于：

\begin{itemize}
\tightlist
\item
  \texttt{ai\_component\_alias.normalized\_alias} 唯一约束
\end{itemize}

因此 \texttt{insert\_component\_alias(...)} 应处理重复别名异常。

\hypertarget{ux653bux51fb-ux7ec4ux4ef6ux5f71ux54cdux5e42ux7b49}{%
\subsubsection{24.3
攻击-组件影响幂等}\label{ux653bux51fb-ux7ec4ux4ef6ux5f71ux54cdux5e42ux7b49}}

基于：

\begin{itemize}
\tightlist
\item
  \texttt{uq\_attack\_component\_impact\_dedup}
\end{itemize}

因此建议使用 upsert 逻辑而非裸插入。

\begin{center}\rule{0.5\linewidth}{0.5pt}\end{center}

\hypertarget{ux65e5ux5fd7ux4e0eux53efux89c2ux6d4bux6027}{%
\subsection{25.
日志与可观测性}\label{ux65e5ux5fd7ux4e0eux53efux89c2ux6d4bux6027}}

数据库模块至少应记录：

\begin{itemize}
\tightlist
\item
  SQL 模块名 / repository 名
\item
  trace\_id（若存在）
\item
  task\_id / attack\_id / raw\_id / component\_id
\item
  执行耗时
\item
  行数影响
\item
  异常类型
\end{itemize}

\hypertarget{ux6ce8ux610f}{%
\subsubsection{注意}\label{ux6ce8ux610f}}

不要把完整 payload、敏感 token、长文本摘要直接打进日志。

\begin{center}\rule{0.5\linewidth}{0.5pt}\end{center}

\hypertarget{ux6d4bux8bd5ux7b56ux7565}{%
\subsection{26. 测试策略}\label{ux6d4bux8bd5ux7b56ux7565}}

\hypertarget{ux5355ux5143ux6d4bux8bd5}{%
\subsubsection{单元测试}\label{ux5355ux5143ux6d4bux8bd5}}

\begin{itemize}
\tightlist
\item
  测试模型映射
\item
  测试 alias 规范化
\item
  测试 SQL 过滤条件
\item
  测试异常转换
\end{itemize}

\hypertarget{ux96c6ux6210ux6d4bux8bd5}{%
\subsubsection{集成测试}\label{ux96c6ux6210ux6d4bux8bd5}}

\begin{itemize}
\tightlist
\item
  使用测试库执行真实 SQL
\item
  校验唯一约束、外键、事务回滚行为
\item
  校验视图查询
\end{itemize}

\hypertarget{ux5efaux8baeux91cdux70b9ux6d4bux8bd5ux573aux666f}{%
\subsubsection{建议重点测试场景}\label{ux5efaux8baeux91cdux70b9ux6d4bux8bd5ux573aux666f}}

\begin{enumerate}
\def\labelenumi{\arabic{enumi}.}
\tightlist
\item
  原始记录重复写入
\item
  主 CVSS 切换
\item
  主 taxonomy 替换
\item
  无法解析的 BOM 入 queue
\item
  \texttt{v\_wp12\_attack\_feed} 是否只返回可供 WP1-2 使用的资产
\end{enumerate}

\begin{center}\rule{0.5\linewidth}{0.5pt}\end{center}

\hypertarget{ux4e0e-agentsintel_agents-ux7684ux8fb9ux754c}{%
\subsection{\texorpdfstring{27. 与 \texttt{agents/intel\_agents/}
的边界}{27. 与 agents/intel\_agents/ 的边界}}\label{ux4e0e-agentsintel_agents-ux7684ux8fb9ux754c}}

\hypertarget{agent-ux8d1fux8d23}{%
\subsubsection{Agent 负责}\label{agent-ux8d1fux8d23}}

\begin{itemize}
\tightlist
\item
  调外部源
\item
  调规则 / LLM
\item
  生成标准化结果对象
\item
  决定何时入库
\end{itemize}

\hypertarget{db-ux8d1fux8d23}{%
\subsubsection{\texorpdfstring{\texttt{db/}
负责}{db/ 负责}}\label{db-ux8d1fux8d23}}

\begin{itemize}
\tightlist
\item
  连接
\item
  事务
\item
  SQL
\item
  模型映射
\item
  读写一致性
\item
  视图读取
\end{itemize}

\hypertarget{ux4e00ux4e2aux975eux5e38ux91cdux8981ux7684ux539fux5219}{%
\subsubsection{一个非常重要的原则}\label{ux4e00ux4e2aux975eux5e38ux91cdux8981ux7684ux539fux5219}}

Agent \textbf{不应直接拼表名写 SQL}。

Agent 只能：

\begin{itemize}
\tightlist
\item
  调 service
\item
  或调 repository（只在非常简单的读场景）
\end{itemize}

\begin{center}\rule{0.5\linewidth}{0.5pt}\end{center}

\end{document}
